\section{Week 14: Modeling in Cognitive Science I: Categorization (GCM)
and Decision Making (Drift-Diffusion)}

\begin{quote}
\textbf{Chapter 6: Examples of modeling in different areas of cognitive
science} · Lecture 1 h · Lab 2 h \textbf{Running data:} simulated
classification data for the classic 5-4 category structure, and
simulated choice-RT data from a drift-diffusion process in a
speed/accuracy × easy/hard design.
\end{quote}

\begin{center}\rule{0.5\linewidth}{0.5pt}\end{center}

\subsection{Key terms}

Important terms of this lesson, with their Persian equivalents (repeated
even if they appeared in earlier weeks).

\begin{longtable}[]{@{}
  >{\raggedright\arraybackslash}p{(\linewidth - 4\tabcolsep) * \real{0.3333}}
  >{\raggedright\arraybackslash}p{(\linewidth - 4\tabcolsep) * \real{0.3333}}
  >{\raggedright\arraybackslash}p{(\linewidth - 4\tabcolsep) * \real{0.3333}}@{}}
\caption{\centering Key terms}\tabularnewline
\toprule\noalign{}
\begin{minipage}[b]{\linewidth}\raggedright
English
\end{minipage} & \begin{minipage}[b]{\linewidth}\raggedright
\RL{\CMfa فارسی}
\end{minipage} & \begin{minipage}[b]{\linewidth}\raggedright
\RL{\CMfa توضیح کوتاه}
\end{minipage} \\
\midrule\noalign{}
\endfirsthead
\toprule\noalign{}
\begin{minipage}[b]{\linewidth}\raggedright
English
\end{minipage} & \begin{minipage}[b]{\linewidth}\raggedright
\RL{\CMfa فارسی}
\end{minipage} & \begin{minipage}[b]{\linewidth}\raggedright
\RL{\CMfa توضیح کوتاه}
\end{minipage} \\
\midrule\noalign{}
\endhead
\bottomrule\noalign{}
\endlastfoot
categorization & \RL{\CMfa دسته‌بندی} & \RL{\CMfa قرار دادن محرک در یک دسته} \\
Generalized Context Model (GCM) & \RL{\CMfa مدل بافت تعمیم‌یافته} & \RL{\CMfa مدل}
exemplar \RL{\CMfa دسته‌بندی بر پایه‌ی شباهت} \\
drift rate & \RL{\CMfa نرخ رانش} & \RL{\CMfa سرعت میانگین جمع شدن شواهد؛ نشانه‌ی
کیفیت شواهد} \\
simulation & \RL{\CMfa شبیه‌سازی} & \RL{\CMfa تولید داده‌ی مصنوعی با اجرای مدل} \\
response time (RT) & \RL{\CMfa زمان واکنش} & \RL{\CMfa فاصله‌ی ارائه‌ی محرک تا
پاسخ} \\
drift-diffusion model (DDM) & \RL{\CMfa مدل} drift-diffusion
\RL{\CMfa (رانش--انتشار)} & \RL{\CMfa انباشت پیوسته‌ی شواهد بین دو مرز} \\
speed--accuracy trade-off & \RL{\CMfa بده‌بستان سرعت و دقت} & \RL{\CMfa سریع‌تر پاسخ
دادن به قیمت خطای بیشتر} \\
accuracy & \RL{\CMfa دقت} & \RL{\CMfa نسبت پاسخ‌های درست} \\
exemplar & \RL{\CMfa نمونه} (exemplar) & \RL{\CMfa یک مورد ذخیره‌شده در حافظه از یک
دسته} \\
prototype & \RL{\CMfa پیش‌نمونه} (prototype) & \RL{\CMfa خلاصه یا میانگین یک دسته} \\
fit / fitting & \RL{\CMfa برازش} & \RL{\CMfa تنظیم پارامترها تا مدل با داده جور
شود} \\
attention weight & \RL{\CMfa وزن توجه} & \RL{\CMfa سهم توجه به هر بُعد محرک} \\
sensitivity (d′) & \RL{\CMfa حساسیت} & \RL{\CMfa فاصله‌ی توزیع علامت و نویز} \\
parameter & \RL{\CMfa پارامتر} & \RL{\CMfa کمیت قابل تنظیم مدل که بین افراد یا شرایط
فرق می‌کند} \\
estimate / estimator & \RL{\CMfa برآورد} / \RL{\CMfa برآوردگر} & \RL{\CMfa مقدار به‌دست‌آمده
برای پارامتر} / \RL{\CMfa روش به دست آوردن آن} \\
EZ-diffusion & EZ-diffusion & \RL{\CMfa برآورد ساده‌ی پارامترهای} DDM \RL{\CMfa از سه
آماره} \\
parameter recovery & \RL{\CMfa بازیابی پارامتر} & \RL{\CMfa آیا برازش، پارامترهای
معلوم شبیه‌سازی را برمی‌گرداند} \\
evidence & \RL{\CMfa شواهد} & \RL{\CMfa اطلاعاتی که برای تصمیم جمع می‌شود} \\
stimulus & \RL{\CMfa محرک} & \RL{\CMfa چیزی که به شرکت‌کننده نشان داده می‌شود} \\
boundary separation & \RL{\CMfa فاصله‌ی مرزها} & \RL{\CMfa مقدار شواهد لازم برای
پاسخ؛ نشانه‌ی احتیاط} \\
non-decision time & \RL{\CMfa زمان غیرتصمیمی} & \RL{\CMfa زمان دیدن محرک و اجرای
حرکت پاسخ} \\
similarity & \RL{\CMfa شباهت} & \RL{\CMfa نزدیکی دو محرک در فضای روان‌شناختی} \\
trial & \RL{\CMfa کوشش} (trial) & \RL{\CMfa یک بار انجام تکلیف توسط شرکت‌کننده} \\
embedding & embedding \RL{\CMfa (بردار معنایی)} & \RL{\CMfa برداری که مدل برای یک
کلمه یا متن می‌سازد} \\
\end{longtable}

\begin{center}\rule{0.5\linewidth}{0.5pt}\end{center}

\subsection{Learning outcomes}

By the end of this week, students will be able to:

\begin{enumerate}
\def\labelenumi{\arabic{enumi}.}
\tightlist
\item
  \textbf{Describe} exemplar and prototype theories of categorization
  and \textbf{formalize} them as the GCM and a prototype model.
\item
  \textbf{Fit} the GCM and the prototype model to classification data by
  MLE, \textbf{compare} them, and \textbf{interpret} attention weights
  and sensitivity.
\item
  \textbf{Describe} the drift-diffusion model (DDM) as the continuous
  limit of the Week 1 random walk, and \textbf{map} its parameters onto
  cognitive constructs.
\item
  \textbf{Estimate} DDM parameters with the closed-form EZ-diffusion
  method and \textbf{validate} it by parameter recovery.
\item
  \textbf{Interpret} experimental effects in terms of model parameters
  (``this manipulation changed caution, not evidence quality'').
\end{enumerate}

\begin{center}\rule{0.5\linewidth}{0.5pt}\end{center}

\subsection{Motivation \& big picture}

The first 13 weeks built a toolkit, mostly on two small running examples
(forgetting and SDT). Chapter 6 shows the toolkit at work on models that
have shaped whole research areas. This week covers two of the most
influential.

\textbf{Categorization.} How do we decide that a new animal is a dog?
\emph{Prototype theory} says that we compare it with a summary of the
category, the average dog. \emph{Exemplar theory} says that we compare
it with stored memories of every dog (and every cat) we have seen. Both
sound plausible, and verbal arguments went back and forth for years.
Formal models made the difference: Medin and Schaffer's 5-4 category
structure was designed so that the two theories make \textbf{opposite
predictions} for specific stimuli. Nosofsky's \textbf{Generalized
Context Model (GCM)} turned exemplar theory into a precise, fit-able
model with interpretable attention weights.

\textbf{Decision making.} In Week 1 we built a random walk that
accumulates evidence to a boundary. Its continuous version, the
\textbf{drift-diffusion model}, is arguably the most successful model in
cognitive psychology: it explains accuracy \emph{and} the full RT
distributions of correct and error responses, with parameters for
evidence quality, caution and non-decision time. It is used to study
ageing, ADHD, depression, learning and the effect of drugs. Fitting the
full DDM is computationally heavy, but the \textbf{EZ-diffusion} method
gives good estimates from three summary statistics with closed-form
equations.

\paragraph{Where is this used?}

\begin{itemize}
\tightlist
\item
  \textbf{Categorization models} explain how people learn concepts
  (medical students learning to recognize skin lesions, children
  learning words), and they are used to design better training: which
  examples to show, and in what order.
\item
  \textbf{Exemplar ideas in machine learning}: nearest-neighbor
  classifiers and kernel methods classify a new case by its similarity
  to stored examples, the same principle as the GCM.
\item
  \textbf{The drift-diffusion model} is one of the most used models in
  cognitive neuroscience and clinical research: ageing, ADHD,
  depression, anxiety, effects of drugs and sleep deprivation. It is
  also used in consumer research (how fast people choose between
  products) and in neuroeconomics.
\item
  \textbf{EZ-diffusion} is popular because it needs only three summary
  statistics per person and no optimization, so it can be applied to
  existing datasets quickly.
\end{itemize}

\begin{center}\rule{0.5\linewidth}{0.5pt}\end{center}

\Needspace{11\baselineskip}

\subsection{Lecture notes}

\paragraph{In plain words}

Two famous models this week. \textbf{Categorization:} how do we decide
that a new animal is a dog? One theory says we compare it with an
average ``ideal dog'' (\textbf{prototype}). Another says we compare it
with all the dogs and cats we remember and see which group it resembles
more overall (\textbf{exemplars}). The \textbf{GCM} turns the exemplar
idea into math: measure the distance between the new item and each
stored example (paying more \textbf{attention} to some features than
others), turn distance into similarity (closer = much more similar), and
choose category A with probability ``total similarity to A / total
similarity to all''. A clever set of stimuli (the 5-4 structure) makes
the two theories predict \emph{opposite} things for some items.
\textbf{Decisions:} the \textbf{drift-diffusion model} is the continuous
version of Week 1's random walk. Evidence drifts toward the correct
boundary at an average speed \(v\) (quality of evidence) while being
jiggled by noise; a response is made when a boundary \(a\) (caution) is
reached; a fixed time \(T_{er}\) is added for seeing and pressing.
\textbf{EZ-diffusion} is a shortcut that computes \(v\), \(a\) and
\(T_{er}\) from just three numbers: accuracy, and the mean and variance
of correct RTs.

\textbf{Suggested timing (60 min):} categorization theories and the 5-4
structure 10 · GCM 15 · prototype model and the critical predictions 5 ·
the drift-diffusion model 15 · EZ-diffusion 10 · interpretation 5.

\subsubsection{1. The 5-4 category structure}

Stimuli vary on four binary dimensions (e.g., the shape, color, size and
position of a geometric figure in Medin \& Schaffer's design; facial
features in later versions). Participants learn to classify 9 training
stimuli into categories A and B, then classify 7 new \textbf{transfer}
stimuli without feedback. The commonly used stimulus set is {[}verify
against Medin \& Schaffer, 1978, or Smith \& Minda, 2000{]}:

\begin{longtable}[]{@{}lll@{}}
\caption{\centering 1. The 5-4 category structure}\tabularnewline
\toprule\noalign{}
Category A & Category B & Transfer \\
\midrule\noalign{}
\endfirsthead
\toprule\noalign{}
Category A & Category B & Transfer \\
\midrule\noalign{}
\endhead
\bottomrule\noalign{}
\endlastfoot
A1 1110 & B1 1100 & T1 1001 \\
A2 1010 & B2 0110 & T2 1000 \\
A3 1011 & B3 0001 & T3 1111 \\
A4 1101 & B4 0000 & T4 0010 \\
A5 0111 & & T5 0101 \\
& & T6 0011 \\
& & T7 0100 \\
\end{longtable}

The category A prototype is 1111 and the B prototype is 0000. The
critical comparison: \textbf{A1 (1110)} shares three features with the A
prototype, while \textbf{A2 (1010)} shares only two. A prototype model
therefore predicts that A1 is classified \emph{better} than A2. But A2
is highly similar to two other A exemplars (A1 and A3) and to no B
exemplar, whereas A1 is highly similar to two \textbf{B} exemplars (B1
and B2). An exemplar model can therefore predict that A2 is classified
\emph{at least as well} as A1, and in the empirical data A2 is typically
classified better.

\subsubsection{2. The Generalized Context Model (GCM)}

\textbf{Distance} between stimulus \(i\) and stored exemplar \(e\), with
attention weights \(w_m \ge 0\), \(\sum_m w_m = 1\) (city-block metric
for separable dimensions):

\[
\mathrm{dist}_{ie} = \sum_{m=1}^{4} w_m\, |x_{im} - x_{em}|.
\]

\textbf{Similarity} decays exponentially with distance, with
\textbf{sensitivity} \(c\):

\[
s_{ie} = \exp(-c \cdot \mathrm{dist}_{ie}).
\]

\textbf{Choice rule} (summed similarity, Luce ratio):

\[
P(A \mid i) = \frac{\sum_{e \in A} s_{ie}}{\sum_{e \in A} s_{ie} + \sum_{e \in B} s_{ie}}.
\]

\textbf{Parameters:} \(c\) (how steeply similarity falls with distance;
large \(c\) means only very close exemplars matter) and three free
attention weights (the fourth is \$1 - \$ the others), so \(k = 4\).
(This is the overloaded \(c\) from the Week 1 notation table; this week
it means sensitivity, not an SDT criterion.)

\textbf{Interpretation of attention weights:} \(w_m\) is the share of
attention given to dimension \(m\). Optimal attention emphasizes
dimensions that are diagnostic of the category. In the 5-4 structure,
dimensions 1 and 3 are the most diagnostic.

\textbf{Prototype model:} the same distance, similarity and choice rule,
but each category is represented by its prototype only:
\(P(A\mid i) = s_{i,P_A}/(s_{i,P_A} + s_{i,P_B})\). Also \(k = 4\), so
the comparison is not affected by parameter counting.

\textbf{Likelihood:} if stimulus \(i\) is classified \(n_i\) times and
called ``A'' \(y_i\) times, then
\(y_i \sim \mathrm{Binomial}(n_i, P(A\mid i))\), fitted by MLE with
multi-start (Weeks 3--4). Weights are best estimated through a
\textbf{softmax reparameterization} (Week 3),
\(w = \mathrm{softmax}(\eta_1, \eta_2, \eta_3, 0)\), and
\(c = e^{\gamma}\).

\begin{quote}
\textbf{Advanced:} the full GCM adds a response-bias parameter \(b_A\),
a similarity exponent (Gaussian vs.~exponential similarity), a metric
parameter \(r\) (city-block \(r = 1\) vs.~Euclidean \(r = 2\)) and
sometimes a response-determinism parameter \(\gamma\) (raising summed
similarities to a power). Lee \& Wagenmakers (chapter 17) develop a
Bayesian, hierarchical GCM.
\end{quote}

\paragraph{Reading the math}

\begin{itemize}
\tightlist
\item
  \textbf{Distance} \(\sum_m w_m|x_{im} - x_{em}|\): count the features
  on which the new item \(i\) and the stored example \(e\) differ, but
  weight each feature by its attention \(w_m\). With 4 binary features
  and equal weights 0.25, one mismatching feature gives distance 0.25.
\item
  \textbf{Similarity} \(e^{-c \cdot \mathrm{dist}}\): distance 0 gives
  similarity 1; with \(c = 2\), distance 0.25 gives \(e^{-0.5} = 0.61\),
  distance 1 gives \(e^{-2} = 0.14\). A larger \(c\) makes similarity
  drop faster, so only very close examples matter.
\item
  \textbf{Choice rule:} add the similarities to all A examples and to
  all B examples, and divide. In the worked example, T3 has total
  similarity 2.79 to A and 1.09 to B, so \(P(A) = 2.79/3.89 = 0.72\).
\end{itemize}

\subsubsection{3. The drift-diffusion model (DDM)}

Take the Week 1 random walk and let the step size and step duration
shrink together. Evidence \(X(t)\) then evolves continuously:

\[
dX = v\,dt + s\,dW(t), \qquad X(0) = z,
\]

where \(dW\) is Gaussian white noise (a Wiener process). A response is
made when \(X\) reaches \(a\) (correct) or \(0\) (error);
\(\mathrm{RT} = T_{\text{decision}} + T_{er}\).

\begin{longtable}[]{@{}
  >{\raggedright\arraybackslash}p{(\linewidth - 4\tabcolsep) * \real{0.3333}}
  >{\raggedright\arraybackslash}p{(\linewidth - 4\tabcolsep) * \real{0.3333}}
  >{\raggedright\arraybackslash}p{(\linewidth - 4\tabcolsep) * \real{0.3333}}@{}}
\caption{\centering 3. The drift-diffusion model (DDM)}\tabularnewline
\toprule\noalign{}
\begin{minipage}[b]{\linewidth}\raggedright
Parameter
\end{minipage} & \begin{minipage}[b]{\linewidth}\raggedright
Meaning
\end{minipage} & \begin{minipage}[b]{\linewidth}\raggedright
Typical manipulation that changes it
\end{minipage} \\
\midrule\noalign{}
\endfirsthead
\toprule\noalign{}
\begin{minipage}[b]{\linewidth}\raggedright
Parameter
\end{minipage} & \begin{minipage}[b]{\linewidth}\raggedright
Meaning
\end{minipage} & \begin{minipage}[b]{\linewidth}\raggedright
Typical manipulation that changes it
\end{minipage} \\
\midrule\noalign{}
\endhead
\bottomrule\noalign{}
\endlastfoot
\(v\) (drift rate) & Quality / rate of evidence accumulation & Stimulus
difficulty; perceptual sensitivity; attention \\
\(a\) (boundary separation) & Response caution (speed--accuracy setting)
& Speed vs.~accuracy instructions \\
\(z\) (starting point) & A priori bias toward one response & Base rates,
payoffs \\
\(T_{er}\) (non-decision time; also written \(t_0\)) & Encoding + motor
time & Response modality, stimulus degradation \\
\(s\) (within-trial noise) & Scaling constant & Fixed (by convention
\(s = 0.1\) in EZ, \(s = 1\) elsewhere) \\
\end{longtable}

The full DDM also includes across-trial variability in drift, starting
point and non-decision time. Week 1, Exercise 5 showed that drift
variability produces slow errors.

\textbf{Signature predictions:} right-skewed RT distributions; a
speed--accuracy trade-off when \(a\) changes; easier stimuli giving
responses that are faster \emph{and} more accurate (higher \(v\)); a
change in \(T_{er}\) shifting the whole RT distribution without changing
its shape.

\paragraph{Reading the math}

\(dX = v\,dt + s\,dW(t)\) reads: ``in a tiny time step \(dt\), the
evidence changes by a deterministic push \(v\,dt\) plus a random kick
\(s\,dW\)''. \(dW\) is Gaussian noise whose variance is proportional to
\(dt\) (so over a step of length \(dt\) the kick has SD \(s\sqrt{dt}\)).
The table of parameters is the important part to memorize: \(v\) =
evidence quality, \(a\) = caution, \(z\) = bias, \(T_{er}\) =
non-decision time.

\subsubsection{4. EZ-diffusion}

Wagenmakers, van der Maas and Grasman (2007) derived closed-form
estimates of \(v\), \(a\) and \(T_{er}\) (assuming \(z = a/2\), no
across-trial variability, \(s = 0.1\)) from three statistics: the
proportion correct \(P_c\), the \textbf{variance} of correct RTs
\(V_{RT}\) (in s²) and the \textbf{mean} of correct RTs \(M_{RT}\) (in
s):

\[
L = \operatorname{logit}(P_c) = \ln\frac{P_c}{1 - P_c}, \qquad
x = \frac{L\,(L P_c^2 - L P_c + P_c - 0.5)}{V_{RT}},
\]

\[
v = \operatorname{sign}(P_c - 0.5)\; s\; x^{1/4}, \qquad
a = \frac{s^2 L}{v},
\]

\[
y = -\frac{v a}{s^2}, \qquad
\mathrm{MDT} = \frac{a}{2v}\cdot\frac{1 - e^{y}}{1 + e^{y}}, \qquad
T_{er} = M_{RT} - \mathrm{MDT}.
\]

\textbf{Edge cases:} the formulas fail at \(P_c = 0.5\) (then \(v = 0\))
and \(P_c = 1\) (then \(L = \infty\)). For \(P_c = 1\), a standard fix
is to use \(P_c = 1 - 1/(2N)\), analogous to the extreme-rate
corrections of Week 2.

\textbf{Strengths and limits.} EZ is fast, needs no optimization and
recovers parameters well when its assumptions hold (lab Part C). It
ignores error RTs and the shape of the RT distribution beyond mean and
variance, and it cannot estimate bias or across-trial variabilities.

\paragraph{Reading the math}

The EZ equations look intimidating but are a recipe; nobody needs to
memorize them.

\begin{enumerate}
\def\labelenumi{\arabic{enumi}.}
\tightlist
\item
  \(L = \ln(P_c/(1-P_c))\): the \textbf{logit} of accuracy (``log odds
  of being correct''). \(P_c = 0.8\) gives \(L = \ln 4 = 1.39\).
\item
  \(x\) combines accuracy with the \textbf{variance} of correct RTs: the
  more variable the RTs for a given accuracy, the lower the drift.
\item
  \(v\) and \(a\) then follow from \(L\) and \(x\); \(\mathrm{MDT}\) is
  the predicted \emph{mean decision time}.
\item
  \(T_{er} = M_{RT} - \mathrm{MDT}\): whatever part of the mean RT is
  not decision time is non-decision time.
\end{enumerate}

\paragraph{For an engineer}

\begin{itemize}
\tightlist
\item
  \textbf{GCM} is a similarity-based classifier with an exponential
  kernel and a weighted L1 (city-block) metric, close to a
  \textbf{kernel / Parzen-window classifier}. The attention weights act
  like learned feature weights; \(c\) is the inverse kernel width.
\item
  The \textbf{DDM} is a \textbf{Wiener process with drift} between two
  absorbing boundaries: the classic first-passage-time problem. The
  continuous-time SPRT (Week 1 box) is exactly this process, with \(v\)
  \ensuremath{\propto} signal amplitude and \(s\) the noise level.
\item
  RT distributions are \textbf{first-passage-time distributions}, which
  are right-skewed for the same reason detection delays of a sequential
  detector are.
\item
  \textbf{EZ-diffusion} is a method-of-moments estimator: it matches
  three sample moments (accuracy, mean and variance of correct RTs) to
  their closed-form model expressions, instead of maximizing a
  likelihood.
\end{itemize}

\paragraph{For the instructor}

\begin{itemize}
\tightlist
\item
  \textbf{Let students be the participants.} Before the GCM, show the 9
  training stimuli of the 5-4 structure as simple drawings (e.g.~shapes
  with 4 binary features), teach the categories with feedback, then test
  A1 and A2. The class's own data often show the exemplar pattern.
\item
  \textbf{Draw the DDM as Week 1's picture with a smooth, noisy path.}
  Then change one parameter at a time on the board (wider boundaries,
  steeper drift, later start) and ask what happens to accuracy and RT.
\item
  \textbf{Do not derive EZ.} Present it as a ``calculator'' and spend
  the time on interpretation (lab Part E): which manipulation changed
  which parameter.
\item
  \textbf{Typical questions.}

  \begin{itemize}
  \tightlist
  \item
    \emph{``Prototype or exemplar, who won?''} For small, poorly
    structured categories exemplar models usually fit better; for large,
    well-structured ones prototypes do well. Both ideas survive in
    modern models.
  \item
    \emph{``Why not just compare mean RTs?''} A slower group could be
    more cautious, have worse evidence, or be slower motorically; only
    the DDM separates them.
  \item
    \emph{``Are the parameters of EZ and the full DDM the same?''} They
    estimate the same quantities, but EZ assumes no bias and no
    trial-to-trial variability.
  \end{itemize}
\end{itemize}

\begin{center}\rule{0.5\linewidth}{0.5pt}\end{center}

\Needspace{11\baselineskip}

\subsection{Worked example}

\subsubsection{Part 1: one GCM prediction by hand}

Stimulus T3 (1111), equal attention \(w_m = 0.25\), sensitivity
\(c = 2\).

\begin{longtable}[]{@{}
  >{\raggedright\arraybackslash}p{(\linewidth - 6\tabcolsep) * \real{0.2500}}
  >{\raggedright\arraybackslash}p{(\linewidth - 6\tabcolsep) * \real{0.2500}}
  >{\raggedright\arraybackslash}p{(\linewidth - 6\tabcolsep) * \real{0.2500}}
  >{\raggedright\arraybackslash}p{(\linewidth - 6\tabcolsep) * \real{0.2500}}@{}}
\caption{\centering Part 1: one GCM prediction by hand}\tabularnewline
\toprule\noalign{}
\begin{minipage}[b]{\linewidth}\raggedright
Exemplar
\end{minipage} & \begin{minipage}[b]{\linewidth}\raggedright
Mismatching dimensions
\end{minipage} & \begin{minipage}[b]{\linewidth}\raggedright
Distance
\end{minipage} & \begin{minipage}[b]{\linewidth}\raggedright
Similarity \(e^{-2\,\mathrm{dist}}\)
\end{minipage} \\
\midrule\noalign{}
\endfirsthead
\toprule\noalign{}
\begin{minipage}[b]{\linewidth}\raggedright
Exemplar
\end{minipage} & \begin{minipage}[b]{\linewidth}\raggedright
Mismatching dimensions
\end{minipage} & \begin{minipage}[b]{\linewidth}\raggedright
Distance
\end{minipage} & \begin{minipage}[b]{\linewidth}\raggedright
Similarity \(e^{-2\,\mathrm{dist}}\)
\end{minipage} \\
\midrule\noalign{}
\endhead
\bottomrule\noalign{}
\endlastfoot
A1 1110 & 1 & 0.25 & 0.607 \\
A2 1010 & 2 & 0.50 & 0.368 \\
A3 1011 & 1 & 0.25 & 0.607 \\
A4 1101 & 1 & 0.25 & 0.607 \\
A5 0111 & 1 & 0.25 & 0.607 \\
\textbf{Sum A} & & & \textbf{2.794} \\
B1 1100 & 2 & 0.50 & 0.368 \\
B2 0110 & 2 & 0.50 & 0.368 \\
B3 0001 & 3 & 0.75 & 0.223 \\
B4 0000 & 4 & 1.00 & 0.135 \\
\textbf{Sum B} & & & \textbf{1.094} \\
\end{longtable}

\[
P(A\mid \text{T3}) = \frac{2.794}{2.794 + 1.094} = 0.719.
\]

\textbf{Interpretation.} T3 is the A prototype, yet the exemplar model
predicts only about 72\% ``A'' responses with these parameters: two B
exemplars (B1 and B2) are only two features away. A larger \(c\)
(sharper similarity) would raise this probability, because the nearest
neighbors (four A exemplars at distance 0.25) would dominate.

\subsubsection{Part 2: EZ-diffusion by hand}

A participant has \(P_c = 0.802\), \(V_{RT} = 0.112\) s²,
\(M_{RT} = 0.723\) s.

\begin{enumerate}
\def\labelenumi{\arabic{enumi}.}
\tightlist
\item
  \(L = \ln(0.802/0.198) = 1.399\).
\item
  \(x = 1.399 \times (1.399 \times 0.643 - 1.399 \times 0.802 + 0.802 - 0.5)/0.112 = 0.998\).
\item
  \(v = 0.1 \times 0.998^{1/4} = 0.100\).
\item
  \(a = 0.01 \times 1.399/0.100 = 0.140\).
\item
  \(y = -0.100 \times 0.140/0.01 = -1.399\), so
  \(\mathrm{MDT} = (0.140/0.200) \times (1 - e^{-1.399})/(1 + e^{-1.399}) = 0.423\)
  s.
\item
  \(T_{er} = 0.723 - 0.423 = 0.300\) s.
\end{enumerate}

\textbf{Interpretation.} Of the mean correct RT of 723 ms, about 423 ms
is spent accumulating evidence and 300 ms on encoding and responding.
Drift 0.10 and boundary 0.14 (in EZ units with \(s = 0.1\)) are typical
values for a moderately difficult task.

\begin{center}\rule{0.5\linewidth}{0.5pt}\end{center}

\Needspace{11\baselineskip}

\subsection{Lab (2 hours)}

\subsubsection{Goals}

\begin{enumerate}
\def\labelenumi{\arabic{enumi}.}
\tightlist
\item
  Formulate and evaluate the GCM and the prototype model for the 5-4 structure;
  simulate data from the GCM.
\item
  Fit both models by MLE, compare them, and examine the critical
  A1-vs-A2 prediction.
\item
  Check parameter recovery for the GCM.
\item
  Simulate a drift-diffusion process, recover its parameters with
  EZ-diffusion, and interpret a 2 × 2 experiment.
\end{enumerate}

\textbf{Timing:} Part A 20 min · Part B 25 min · Part C 15 min · Part D
40 min · Part E 20 min.

\subsubsection{Part B: fit and compare (25 min)}

\textbf{Interpret:} 1. Which model is preferred, and how strong is the
evidence? (\(\Delta\mathrm{AIC} = 3.3\), so Akaike weights of about 0.84
vs.~0.16: a moderate preference, not a decisive one.) 2. Look at A1 and
A2. The data show A2 (0.90) classified better than A1 (0.75). The GCM
predicts this ordering (0.88 vs.~0.84); the prototype model cannot (0.82
vs.~0.82). Why does this comparison discriminate between the theories
more sharply than the overall fit statistic? 3. Interpret the GCM
attention weights. Which dimensions did the participant attend to, and
does that make sense given the category structure?

\subsubsection{Part C: is the GCM recoverable at this design? (15 min)}

\textbf{What to notice:} the attention weights are recovered well (r =
0.76--0.91), but the sensitivity \(c\) is essentially \textbf{not
recoverable} at this design (r = 0.08). With only 16 stimuli and 20
trials each, a large \(c\) and a moderate \(c\) can produce very similar
response proportions (once \(c\) is large, predictions saturate), and
\(c\) trades off with how concentrated the weights are. So we should
interpret the attention weights from Part B but \textbf{not} the value
\(\hat c = 3.82\). More trials, a wider range of stimuli, or a Bayesian
prior on \(c\) would help. This is exactly the kind of result a recovery
study is for (Week 7).

\subsubsection{Part E: interpret the 2 × 2 experiment (20 min)}

Which EZ parameter changes with the \textbf{instruction} manipulation,
and which with \textbf{difficulty}? Did either manipulation affect the
parameter it should \emph{not} affect? 2. Describe how the RT
distribution changes from speed to accuracy instructions (Part D
figure). Which part of the distribution changes most: the leading edge
or the tail? 3. A clinical study finds that patients are slower than
controls but equally accurate. Sketch two different DDM accounts of this
result, and say which data would distinguish them.

\subsubsection{Exercises (increasing difficulty)}

\textbf{Exercise 1 \ensuremath{\star}.} Compute \(P(A\mid\text{A1})\) and
\(P(A\mid\text{A2})\) by hand for equal weights and \(c = 2\) under both
models. Do they reproduce the critical prediction?

\textbf{Exercise 2 \ensuremath{\star}.} Apply EZ to the worked example's statistics and
to \(P_c = 1.0\) with \(N = 100\) trials (use the edge correction). What
happens without the correction?

\textbf{Exercise 3 \ensuremath{\star}\ensuremath{\star} (prototype data).} Generate data from the
\textbf{prototype} model and fit both models. Does the model comparison
identify the generating model? Run a small model recovery (20 datasets
per generating model) and build the 2 × 2 confusion matrix.

\textbf{Exercise 4 \ensuremath{\star}\ensuremath{\star} (EZ recovery and bias).} For 200 simulated
participants with random \((v, a, T_{er})\) and only 100 trials each,
compare EZ estimates with the true values. Which parameter is recovered
worst? Then add across-trial drift variability to the simulator (Week 1,
Exercise 5) and check which EZ estimate becomes biased.

\(c \sim \mathrm{Gamma}(2, \text{scale } 2)\) ·
\(w \sim \mathrm{Dirichlet}(1, 1, 1, 1)\) Compare the posterior of the
attention weights with the MLE, and run a posterior predictive check on
the A1-vs-A2 difference. (See Lee \& Wagenmakers, chapter 17, for a
hierarchical version.)

\subsubsection{Lab check-off}

Show: (1) the fit table and prediction plot for the GCM vs.~the
prototype model with your answer to question 2, (2) the EZ table for the
2 × 2 design, and (3) your answer to Part E, question 1.

\begin{center}\rule{0.5\linewidth}{0.5pt}\end{center}

\subsection{Connection to AI \& LLMs}

\textbf{Exemplar models in a learned ``semantic space''.} The GCM needs
a psychological space in which to measure distances. Traditionally, that
space comes from similarity ratings and multidimensional scaling.
Language models offer an alternative: every word or sentence is mapped
to an \textbf{embedding vector}, and the distances between vectors are
often used as a stand-in for psychological similarity. Combining
embeddings with the GCM gives a model that can classify \emph{any} new
word with no extra similarity ratings. Whether embedding distances match
human similarity is an empirical question, and you can test it with the
tools of this course.

How do the three predictions change, and why does ``nutcracker'' change
the most? - \ensuremath{\star}\ensuremath{\star} Suppose 20 participants classified the three probes, with
19, 1 and 12 ``fruit'' responses. Fit \(c\) by MLE (binomial likelihood,
Week 4) and run a \(G^2\) test. Does the embedding space with a single
parameter fit? - \ensuremath{\star}\ensuremath{\star} Discussion: the embeddings were learned from text,
not from perception. For which categories might text-based similarity
differ systematically from human similarity (think of color or shape)?

\begin{center}\rule{0.5\linewidth}{0.5pt}\end{center}

\subsection{Common mistakes \& misconceptions}

\begin{enumerate}
\def\labelenumi{\arabic{enumi}.}
\tightlist
\item
  \textbf{Unconstrained attention weights.} Without a sum-to-one
  constraint, \(c\) and \(w\) trade off perfectly (non-identifiable).
  Use a softmax reparameterization.
\item
  \textbf{Judging exemplar vs.~prototype models by overall fit only.}
  The theoretically diagnostic predictions (A1 vs.~A2) are often more
  informative than small differences in AIC.
\item
  \textbf{Averaging classification data over participants} who attend to
  different dimensions. This can make an exemplar-like individual look
  prototype-like (and vice versa); recall Week 5's aggregation
  artifacts.
\item
  \textbf{Using EZ with incorrect units.} RTs must be in
  \textbf{seconds} and the variance in s², with \(s = 0.1\). Mixing ms
  and s gives nonsense.
\item
  \textbf{Applying EZ when its assumptions are violated} (strong bias,
  slow or fast errors). Check error RTs; use the full DDM if needed.
\item
  \textbf{``A slower group has a lower drift rate.''} Slowness can come
  from a higher boundary, a lower drift, or a longer non-decision time.
  Only a model fitted to accuracy \emph{and} RT distributions can tell
  which.
\item
  \textbf{Forgetting the edge cases} \(P_c = 1\) or \(P_c \le 0.5\) in
  EZ.
\end{enumerate}

\begin{center}\rule{0.5\linewidth}{0.5pt}\end{center}

\subsection{Check your understanding}

\textbf{1.} State the core difference between exemplar and prototype
theories in one sentence each. \textbf{Answer} Exemplar: a category is
represented by all stored instances, and classification depends on
summed similarity to them. Prototype: a category is represented by a
single summary (central tendency), and classification depends on
similarity to that summary.

\textbf{2.} In the 5-4 structure, why is the A1-vs-A2 comparison
diagnostic? \textbf{Answer} A1 is closer to the A prototype (3 of 4
features), so a prototype model predicts it is classified better than
A2. But A2 is very similar to other A exemplars and to no B exemplar,
while A1 is similar to two B exemplars, so the exemplar model can
predict A2 \ensuremath{\geq} A1. The theories make opposite ordinal predictions.

\textbf{3.} What do the GCM parameters \(c\) and \(w_m\) represent?
\textbf{Answer} \(c\) is sensitivity: how quickly similarity decreases
with psychological distance (large \(c\) means only very close exemplars
matter). \(w_m\) is the proportion of attention allocated to dimension
\(m\), which stretches distances along attended dimensions.

\textbf{4.} Which DDM parameter would you expect to change under (a)
speed vs.~accuracy instructions, (b) harder stimuli, (c) a slower motor
response (e.g., foot instead of finger)? \textbf{Answer} (a) Boundary
separation \(a\). (b) Drift rate \(v\). (c) Non-decision time
\(T_{er}\).

\textbf{5.} What three statistics does EZ-diffusion need, and what are
its main assumptions? \textbf{Answer} Proportion correct, the variance
of correct RTs, and the mean of correct RTs. It assumes an unbiased
starting point (\(z = a/2\)), no across-trial variability in parameters,
and the scaling \(s = 0.1\).

\textbf{6.} How is the DDM related to the random walk model of Week 1?
\textbf{Answer} It is the continuous-time, continuous-evidence limit of
the random walk: as step size and step duration go to zero together, the
random walk converges to a Wiener diffusion process with drift \(v\) and
noise \(s\).

\begin{center}\rule{0.5\linewidth}{0.5pt}\end{center}

\subsection{Assignment}

\begin{itemize}
\tightlist
\item
  \textbf{A5 (MCMC from scratch)} is due at the start of this week's
  lecture.
\end{itemize}

\begin{center}\rule{0.5\linewidth}{0.5pt}\end{center}

\subsection{Further reading}

\begin{itemize}
\tightlist
\item
  \textbf{Lee \& Wagenmakers (2014), chapter 17.} The Generalized
  Context Model, fitted in a Bayesian (and hierarchical) way to
  category-learning data. \emph{(Main reading for categorization.)}
\item
  \textbf{Farrell \& Lewandowsky (2018)}: the chapter on models of
  choice response time {[}verify chapter number{]} for the
  drift-diffusion model and its relatives.
\item
  \textbf{Verguts (2022)}: the chapters on categorization and on
  decision making {[}verify chapters{]}.
\item
  \emph{Optional, outside the course list:} Wagenmakers, E.-J., van der
  Maas, H. L. J., \& Grasman, R. P. P. P. (2007). An EZ-diffusion model
  for response time and accuracy. \emph{Psychonomic Bulletin \& Review,
  14}, 3--22. Medin, D. L., \& Schaffer, M. M. (1978). Context theory of
  classification learning. \emph{Psychological Review, 85}, 207--238.
\end{itemize}

\emph{Preview for Week 15:} learning from rewards (reinforcement
learning), the mind as a Bayesian observer, and how cognitive models and
AI systems inform each other.
